% Hodina 40, 1. ročník SPŠ. Rozsah: Rozpis_vyuky_MAT_1_rocnik.tex, řádek 40.

% Plán: 1_rocnik/tematicke_plany/SPS.pdf; společná projektová šablona.

% Učebnice: ucebnice/01_Zakladni_poznatky/Mocniny_a_odmocniny/Mocniny a Odmocniny.tex

% Původ úloh: Vlastní zadání a vlastní formulace výkladu; metodická inspirace zdroji.

% Metodický zdroj: externi_zdroje/realisticky/Mocniny/10_Mocniny_-_shrnuti.pdf

% URL: http://www.realisticky.cz/ucebnice/01%20Matematika%20S%C5%A0/01%20Z%C3%A1kladn%C3%AD%20poznatky/06%20Mocniny/10%20Mocniny%20-%20shrnut%C3%AD.pdf

% Metodický zdroj: externi_zdroje/realisticky/Odmocniny/09_Usmernovani_zlomku.pdf

% URL: http://www.realisticky.cz/ucebnice/01%20Matematika%20S%C5%A0/01%20Z%C3%A1kladn%C3%AD%20poznatky/02%20%C4%8C%C3%ADseln%C3%A9%20obory/09%20Usm%C4%9Br%C5%88ov%C3%A1n%C3%AD%20zlomk%C5%AF.pdf

\documentclass[11pt,a4paper]{article}

% Spolecny styl pro ucitelske pripravy MAT (Work/MAT)
% Verze 1.1, 2026-10-23
\NeedsTeXFormat{LaTeX2e}
\ProvidesPackage{priprava}[2026/10/23 v1.1 Sablona ucitelskych priprav MAT]

% Zakladni typografie a cestina
\RequirePackage[T1]{fontenc}
\RequirePackage[utf8]{inputenc}
\RequirePackage[czech]{babel}
\RequirePackage{lmodern}
\RequirePackage{microtype}
\RequirePackage{amsmath,amssymb}
\RequirePackage{xcolor}
\RequirePackage{enumitem}
\RequirePackage[most]{tcolorbox}
\RequirePackage{tabularx,booktabs,array}
\RequirePackage{fancyhdr}
\RequirePackage{needspace}
\RequirePackage{geometry}
\RequirePackage[hidelinks]{hyperref}

\geometry{a4paper,top=17mm,bottom=18mm,left=23mm,right=23mm}
\definecolor{notegray}{gray}{0.39}
\setlength{\parindent}{0pt}
\setlength{\parskip}{2.5pt plus 1pt minus 0.5pt}
\setlist[itemize]{topsep=2pt,itemsep=1pt,parsep=0pt,leftmargin=1.45em}
\setlist[enumerate]{topsep=2pt,itemsep=2pt,parsep=0pt,leftmargin=1.65em}
\setcounter{secnumdepth}{0}
\renewcommand{\arraystretch}{1.12}

% Zahlavi: vlevo rocnik | globalni tema, vpravo cislo: nazev hodiny
% Zapati zustava zcela prazdne.
\pagestyle{fancy}
\fancyhf{}
\renewcommand{\headrulewidth}{0.35pt}
\renewcommand{\footrulewidth}{0pt}
\setlength{\headheight}{14pt}
\newcommand{\nastavPripravu}[4]{%
  \fancyhead[L]{\small #1\enspace\textbar\enspace #2}%
  \fancyhead[R]{\small #3: #4}%
}

% Minimalisticke cernobile boxy. Kazda definice zvlast.
\tcbset{commonlesson/.style={enhanced,breakable=false,colback=white,colframe=black,
  colbacktitle=white,coltitle=black,boxrule=0.45pt,
  arc=0.8mm,outer arc=0.8mm,left=3mm,right=3mm,top=1.6mm,bottom=1.6mm,
  toptitle=1mm,bottomtitle=0.7mm,fonttitle=\bfseries,
  attach boxed title to top left={xshift=2.5mm,yshift=-1.15mm},
  boxed title style={colback=white,colframe=white,boxrule=0pt,left=1mm,right=1mm}}}
\newtcolorbox{zadaniBox}[1][]{commonlesson,title={Zadání},#1}
\newtcolorbox{prikladBox}[1][]{commonlesson,title={Příklad},#1}
\newtcolorbox{definiceBox}[1][]{commonlesson,title={Definice},#1}

% Spolecny box pro zadani a jeho reseni, bez mezery mezi castmi.
% Pouziti:
%   \begin{prikladReseniBox}
%     Zadani ulohy.
%     \tcblower
%     \begin{align*} ... \end{align*}
%   \end{prikladReseniBox}
% Spodni cast ma jemny sedy podklad a delici linku; lze v ni sazet
% vice radku matematiky. Obe casti maji jen jeden spolecny ramecek.
\tcbset{solutionlesson/.style={
  commonlesson,
  enhanced jigsaw,
  breakable,
  colbacklower=black!3,
  segmentation style={solid,draw=black!25,line width=0.35pt},
  fontlower=\small\color{notegray},
  before lower={\noindent\textbf{Řešení:}\enspace},
  leftlower=3mm,rightlower=3mm
}}
\newtcolorbox{prikladReseniBox}[1][]{solutionlesson,title={Příklad},#1}
\newtcolorbox{zadaniReseniBox}[1][]{solutionlesson,title={Zadání},#1}

% Pomucky pro autora pripravy: sedive poznamky; \reseni pro starsi pripravy
\newcommand{\poznamka}[1]{\par\begingroup\small\color{notegray}\itshape
  \textnormal{\textbf{Poznámka pro učitele.}}\enspace #1\par\endgroup}
\newcommand{\reseni}[1]{\par\begingroup\small\color{notegray}
  \textbf{Očekávané řešení:}\enspace #1\par\endgroup}
\newcommand{\blok}[1]{\needspace{5\baselineskip}\subsection*{#1}}
\newcommand{\vyrok}[1]{\uv{#1}}


\nastavPripravu{1. ročník}{Odmocniny}{40}{Průběžný test a rozbor}

\begin{document}

\blok{Výklad a souvislosti}

Rozpis vyhrazuje 15–20 minut průběžnému testu z mocnin a odmocnin. Zbytek hodiny slouží krátkému rozboru a navazujícímu procvičování. Zde je učitelská metodika a zásoba k nápravě; konkrétní test, varianty a jejich schválení řeší samostatně skill \texttt{vytvor-test}.

\begin{definiceBox}[title={Připomenutí definice},breakable=false]
\textbf{Definiční obor} je soubor všech přípustných hodnot proměnných v původním výrazu; sledujeme jej i tehdy, když se výraz po úpravě zkrátí.
\end{definiceBox}


Před testem připrav schválené zadání, stanov způsob zápisu postupů a sděl organizační pravidla. Tato příprava nepovoluje žádné další pomůcky. Během písemné části nedoplňuj matematickou nápovědu. Po odevzdání porovnej několik rozhodujících kroků, například určení základu, práci se záporným exponentem a kontrolu radikandu. Teprve potom zvol cílené procvičení podle zjištěných potíží. Nevyžaduj projít všechny rezervní položky za zbytek hodiny.

\blok{Řešené příklady}

\begin{prikladReseniBox}[title={Příklad 1},breakable=false]
Pro následný rozbor vysvětli, proč $\sqrt{(-4)^2}$ není $-4$.
\tcblower
Druhá odmocnina je nezáporná: $\sqrt{16}=4$. Obecně platí $\sqrt{a^2}=|a|$.
\end{prikladReseniBox}


\begin{prikladReseniBox}[title={Příklad 2},breakable=false]
Při rozboru uprav $a^{-2}:a^{-3}$.
\tcblower
Oba původní základy musí být nenulové, tedy $a\ne0$. Výsledek je $a^{-2-(-3)}=a$. Podmínku zachováme.
\end{prikladReseniBox}


\blok{Samostatné procvičení}

\begin{zadaniReseniBox}[title={Procvičení},breakable=false]
\begin{enumerate}[label=\alph*)]
\item Jako cílenou nápravu vypočítej $(-2)^{-3}$.
\item Při procvičování uprav $\sqrt{32}+\sqrt8$.
\item Při nápravě usměrni $\frac2{\sqrt3+1}$.
\end{enumerate}
\tcblower
\begin{enumerate}[label=\alph*)]
\item $-1/8$, protože převrácená hodnota $(-2)^3=-8$ je záporná.
\item $4\sqrt2+2\sqrt2=6\sqrt2$.
\item $\sqrt3-1$, protože po rozšíření vznikne jmenovatel $3-1=2$.
\end{enumerate}
\end{zadaniReseniBox}


\poznamka{Nejde o žákovské testové zadání. Pro následné procvičení vyber nejvýše potřebné položky. Uznávání postupů a případné bodování stanoví učitel u konkrétního testu; zde nejsou body ani známkové hranice. Poslední krátké ověření ať žák provede samostatně.}

\end{document}
